% !TeX program = xelatex
% !TeX spellcheck = de_DE
\documentclass{beamer}
\usepackage{fontspec}
\usepackage{polyglossia}
\usepackage{hyperref}
\usepackage{xspace}
\usepackage{array}
\usepackage{fontawesome}
\usepackage[normalem]{ulem}

%\setsansfont{Calibri} kann ich sonst nicht kompilieren
\setdefaultlanguage{german}
\usetheme{UHH}

\def\authors#1#2{\author[#1]{#2}}
\def\email#1{\texttt{#1}}
\title{Omega Automaten und LTL}
\subtitle{Seminararbeit zu der Veranstaltung \\Formale Grundlagen der Informatik 3}
\authors{Jennifer~Soltau \& Felix~Ortmann}{Jennifer~Soltau und Felix~Ortmann\\ \email{\{2soltau,0ortmann\}@informatik.uni-hamburg.de}}
\institute{Universität Hamburg\\Fachbereich Informatik\\Vogt-Kölln-Straße 30\\22527 Hamburg}


\date{22-26. Februar 2016}
\titlegraphic{\includegraphics[width=3cm]{img/uhh.pdf}}

\setcounter{tocdepth}{1}
\def\tocname{Gliederung des Vortrags}
\AtBeginSection{\frame{\frametitle{\tocname}\tableofcontents[currentsection]}}

\newcolumntype{x}[1]{>{\centering\arraybackslash}m{#1}}
\setbeamertemplate{headline}[default]

\defbeamertemplate*{title page}{customized}[1][]
{
	{\color{leuchtrot}
	\usebeamerfont{title}\centering\inserttitle\par
	\centering\usebeamerfont{subtitle}\insertsubtitle\par}
	\bigskip
	\centering\usebeamerfont{author}\insertauthor\par
	\bigskip
	\usebeamerfont{institute}\insertinstitute\par
	\bigskip
	\usebeamerfont{date}\insertdate\par
}

\begin{document}

\fontsize{14pt}{14pt}
	
% Titelfolie
{\setbeamercolor{title page}{bg=leuchtrot}
\begin{frame}%
	\titlepage
\end{frame}}

% Gliederung
\begin{frame}{\tocname}
	\tableofcontents
\end{frame}

\section{Einführung}
\section{Omega Automaten}
% Endliche Automaten

\begin{frame}{Endliche Automaten}
\begin{block}{Definition Endlicher Automat}
Ein endlicher Automat A ist ein Tupel ($\Sigma, S, S^0, \rho, F$) mit
\vspace{1em}
\begin{itemize}
\setlength\itemsep{1em}
\item $\Sigma$ ist nicht leeres Alphabet,
\item $S$ ist endliche nicht leere Menge von Zuständen,
\item $S^0\subseteq S$ ist nicht leere Menge von Startzuständen,
\item $F\subseteq S$ ist Menge von Endzuständen und
\item $\rho : S \times \Sigma \rightarrow 2^S$ Transitionsfunktion.
\end{itemize}
\end{block}
\end{frame}

\begin{frame}{Endlicher Automat}
\begin{block}{Definition Endlicher Automat (Fortsetzung)}
\begin{itemize}
\setlength\itemsep{1em}
\item Ist $|S^0|=1$ und $|\rho(s,a)|\leq 1$ für alle $s\in S$ und $a\in \Sigma$, dann ist der Automat deterministisch.
\pause
\item Ein Lauf $r$ auf dem Wort $w=a_0,...,a_{n-1}\in \Sigma^\ast$ ist eine Sequenz von Zuständen $s_0,...,s_n$, sodass $s_0\in S^0$ und $s_{i+1}\in \rho(s_i,a_i)$ für $0\leq i < n$ gilt.
\item Ein Lauf $r$ ist akzeptierend, wenn $s_n\in F$ gilt.
\end{itemize}
\end{block}
\end{frame}

% Alternierende endliche Automaten

\begin{frame}{Alternierende Endliche Automaten}
\begin{block}{Alternation}
\begin{itemize}[<+->]
\setlength\itemsep{1em}
\item Nichtdeterminismus ermöglicht existenzielle Entscheidung
\item Duales Konzept ist universelle Wahl
\item $\rightarrow$ Alternation verbindet beide Konzepte
\end{itemize}
\end{block}
\end{frame}

\begin{frame}{Alternierende Endliche Automaten}
\begin{block}{Die Menge $\mathcal{B}^+(X)$}
\begin{itemize}
\setlength\itemsep{1em}
\item $\mathcal{B}^+(X)$ ist die Menge der positiven boolschen Formeln über der Menge X.
\item Beinhaltet auch die Formeln $\top$ (true) und $\bot$ (false).
\pause
\item $Y \subseteq X$ erfüllt eine Formel $\theta \in \mathcal{B}^+(X)$, wenn die Wahrheitsbelegung, die $true$ für die Elemente von $Y$ und $false$ für die Elemente von $X-Y$ zuordnet, $\theta$ erfüllt.
\pause
\item Beispiel \cite{vardi+96}: $Y=\{s_1,s_3\}$ erfüllt die Formel $\theta = (s_1 \vee s_2)\wedge (s_3\vee s_4)$\\
\end{itemize}
\end{block}
\end{frame}

\begin{frame}{Alternierende Endliche Automaten}
\begin{block}{Definition Alternierender Endlicher Automat}
\begin{itemize}
\setlength\itemsep{1em}
\item Unterscheidet sich nur in der Transitionsfunktion von einem nichtdeterministischen Automaten
\item Formal: $\rho: S\times \Sigma \rightarrow \mathcal{B}^+(S)$
\end{itemize}
\end{block}
\end{frame}

\begin{frame}{Alternierende Endliche Automaten}
\begin{block}{Definition Alternierender Endlicher Automat (Fortsetzung)}
Ein Lauf r auf einem $w=a_0,a_1,...,a_{n-1}$ ist ein (unendlicher) S-beschrifteter Baum r mit
\begin{itemize}
\setlength\itemsep{1em}
\item $r(\epsilon)=s^0$,
\pause
\item und für die folgenden Knoten gilt: wenn $|x|=i, r(x)=s$ und $\rho(s,a_i)=\theta$, dann hat $x$ $k$ Kinder $x_1,...,x_k$ für ein $k \leq |S|$ und $r(x_1),...,r(x_k)$ erfüllt $\theta$.
\end{itemize}
\pause
\vspace{1em}
Ein Lauf r wird akzeptiert, wenn alle Knoten in der Tiefe $n$ mit Zuständen aus $F$ markiert sind.
\end{block}
\end{frame}

% Büchi-Automaten

\begin{frame}{Büchi-Automaten}
\begin{block}{Definition Büchi-Automat}
\begin{itemize}
\setlength\itemsep{1em}
\item Büchi-Automaten akzeptieren unendliche Wörter $a_1,a_2,... \in \Sigma^\omega$
\item Definition wie endlicher Automat, aber besitzen nur einen Startzustand $s^0$
\pause
\item Ein Lauf $r$ auf dem Wort $w=a_0,a_1,...\in \Sigma^\omega$ ist eine unendliche Folge von Zuständen $s_0,s_1,...$, wobei $s_0\in S^0$ und $s_{i+1}\in \rho(s_i,a_i)$ für alle $i\geq0$ gilt.
\pause
\item $lim(r)$ ist Menge \{$s|s=s_i$ für unendlich viele i's\}
\item Ein Lauf $r$ ist akzeptierend, wenn $lim(r)\cap F \neq \emptyset$ gilt.
\end{itemize}
%Ein Büchi-Automat A ist ein Tupel ($\Sigma, S, s^0, \rho, F$) mit
%\begin{itemize}
%\item $\Sigma$ ist nicht leeres Alphabet,
%\item $S$ ist endliche nicht leere Menge von Zuständen,
%\item $s^0$ ist ein Startzustand,
%\item $F\subseteq S$ ist Menge von Endzuständen und
%\item $\rho : S \times \Sigma \rightarrow 2^S$ Transitionsfunktion.
%\end{itemize}
\end{block}
\end{frame}

\begin{frame}{Alternierende Büchi-Automaten}
\begin{block}{Definition Alternierender Büchi-Automat}
\begin{itemize}
\setlength\itemsep{1em}
\item Unterschied zu Büchi-Automaten liegt nur in der Definition der Transaktionsfunktion
\item Formal: $\rho: S\times \Sigma \rightarrow \mathcal{B}^+(S)$
\end{itemize}
%Ein Büchi-Automat A ist ein Tupel ($\Sigma, S, s^0, \rho, F$) mit
%\begin{itemize}
%\item $\Sigma$ ist nicht leeres Alphabet,
%\item $S$ ist endliche nicht leere Menge von Zuständen,
%\item $s^0$ ist ein Startzustand,
%\item $F\subseteq S$ ist Menge von Endzuständen und
%\item \textbf{$\rho : S \times \Sigma \rightarrow \mathcal{B}^+(S)$} Transitionsfunktion.
%\end{itemize}
\end{block}
\end{frame}

\begin{frame}{Alternierende Büchi-Automaten}
\begin{block}{Definition Alternierender Büchi-Automat (Fortsetzung)}
Ein Lauf r auf einem unendlichen Wort $w=a_0,a_1,...$ ist ein (unendlicher) S-beschrifteter Baum r mit
\begin{itemize}
\setlength\itemsep{1em}
\item $r(\epsilon)=s^0$,
\pause
\item und für die folgenden Knoten gilt: wenn $|x|=i, r(x)=s$ und $\rho(s,a_i)=\theta$, dann hat $x$ $k$ Kinder $x_1,...,x_k$ für ein $k \leq |S|$ und $r(x_1),...,r(x_k)$ erfüllt $\theta$.
\end{itemize}
\pause
\vspace{1em}
Ein Lauf r wird akzeptiert, wenn jeder unendliche Zweig unendlich viele Markierungen aus der Endzustandsmenge enthält.
\end{block}
\end{frame}


\section{LTL}
\section{Alternierende Büchi Automaten und LTL}
\section{LTL auf normalen Büchi Automaten}
\begin{frame}{LTL auf Büchi-Automaten}
\begin{block}{Schritte:}
\begin{itemize}
\item Überführung von alternierenden Büchi-Automaten zu normalen nichtdeterministischen Büchi-Automaten
\item Schlussfolgerung für LTL mit Büchi-Automaten
\end{itemize}
\end{block}
\end{frame}

\begin{frame}{Alternierende Büchi-Automaten zu Büchi-Automaten}
\begin{block}{Satz}
Sei A ein alternierender Büchi-Automat. Dann gibt es einen nicht nichtdeterministischen Büchi-Automaten $A_n$ sodass $L_{\omega}(A_n)=L_{\omega}(A)$ gilt.
\end{block}
\end{frame}

\begin{frame}{Alternierende Büchi-Automaten zu Büchi-Automaten}
\begin{block}{Beweis}
\begin{itemize}
\item $A_n$ nimmt Lauf von A an und merkt sich gesamte Ebene eines Laufes
\item einzelnen Ebenen eines Laufes sind Bäume
\item $A_n$ merkt sich das Auftreten von akzeptierenden Zuständen
\item Aufteilung der Ebene eines Laufs in zwei Mengen
\begin{itemize}
\item Zweige, die vor kurzem einen Endzustand erreicht haben
\item Zweige, die vor kurzem keinen Endzustand erreicht haben
\end{itemize}
\end{itemize}
\end{block}
\end{frame}

\begin{frame}{Alternierende Büchi-Automaten zu Büchi-Automaten}
\begin{block}{Beweis (Konstruktion)}
Gegeben $A=(\Sigma,S, s^0, \rho, F)$. Dann $A_n=(\Sigma,S_n,S^0,\rho_n,F_n)$ mit
\begin{itemize}
\item $S_n=2^S\times 2^S$
\item $F_n=\emptyset\times 2^S$
\item und für die Transitionsfunktion gilt: $\rightarrow$ nächste Folie
\end{itemize}
\end{block}
\end{frame}

\begin{frame}{Alternierende Büchi-Automaten zu Büchi-Automaten}
\begin{block}{Beweis (Transitionsfunktion Teil 1)}
Für die Transitionsfunktion gilt:
\begin{itemize}
\item für $U\neq\emptyset$:
\begin{align*}
\rho_n((U,V),a)=\{(U',V')| \text{es gibt} X,Y\subseteq S\text{ sodass}&\\
X \text{ erfüllt} \bigwedge_{t \in U}\rho(t,a),\\
Y \text{ erfüllt} \bigwedge_{t_\in V}\rho(t,a),\\
U'=X-F, \text{ und } V'=Y\cup(X\cap F)\},
\end{align*}
\end{itemize}
\end{block}
\end{frame}

\begin{frame}{Alternierende Büchi-Automaten zu Büchi-Automaten}
\begin{block}{Beweis (Transitionsfunktion Teil 2)}
Für die Transitionsfunktion gilt:
\begin{itemize}
\item für $U=\emptyset$:
\begin{align*}
\rho_n((\emptyset,V),a)=\{(U',V')|\text{ es gibt } X,Y\subseteq S \text{ sodass }\\
Y \text{ erfüllt } \bigwedge_{t_\in V}\rho(t,a),\\
U'=Y-F\text{, und } V'=Y\cap F\}.
\end{align*}
\end{itemize}
\end{block}
\end{frame}

\begin{frame}{LTL und Büchi-Automaten}
\begin{block}{Korollar \cite{vardi+96,vardi+94}}
Gegeben eine LTL Formel $\varphi$ ist es möglich einen Büchi-Automaten $A_{\varphi}=(\Sigma,S, S^0,\rho,F)$ zu konstruieren, wobei $\Sigma=2^{Prop}$ und $|S|$ in $2^{O(|\varphi|)}$, sodass $L_{\omega}(A_n)$ genau die Menge der Berechnungen, die die Formel $\varphi$ erfüllen.
\end{block}
\end{frame} 
\section{Ausblick: Modelchecking}
\section{Zusammenfassung}

% Literatur
\section{\bibname}
\begin{frame}[allowframebreaks]{\bibname}
	\AtBeginSection{}
	\nocite{*}
	\bibliographystyle{abbrv}
	\bibliography{bib}
\end{frame}
\begin{frame}[allowframebreaks]{Bildquellen}
\end{frame}


\end{document}